\begin{table}[H]
	\centering
	\caption{Calibration fraction sensitivity in small samples (500 replications, $\alpha = 0.01$). The forecaster assumes Normal innovations under misspecified heavy-tailed DGPs. Green\% denotes the percentage of replications satisfying the Basel Traffic Light Green zone (scaled to sample size).}
	\label{tab:fc_small_sample}
	\footnotesize
	\begin{tabular}{@{}llcccc@{}}
		\toprule
		DGP & $(T, f_c)$ & Mean $\qV$ & Std $\qV$ & $\hat\pi$ & Green\% \\
		\midrule
		\multicolumn{6}{@{}l}{\textit{Student-$t(5)$}} \\[2pt]
		& (250, 0.40) & 0.0120 & 0.0179 & 0.013 & 62 \\
		& (250, 0.50) & 0.0144 & 0.0162 & 0.011 & 75 \\
		& (250, 0.60) & 0.0167 & 0.0178 & 0.010 & 81 \\
		& (250, 0.70) & 0.0190 & 0.0203 & 0.007 & 72 \\
		& (250, 0.80) & 0.0215 & 0.0251 & 0.007 & 79 \\
		\addlinespace
		& (500, 0.40) & 0.0201 & 0.0188 & 0.007 & 85 \\
		& (500, 0.50) & 0.0111 & 0.0112 & 0.009 & 83 \\
		& (500, 0.60) & 0.0132 & 0.0135 & 0.010 & 76 \\
		& (500, 0.70) & 0.0143 & 0.0121 & 0.007 & 78 \\
		& (500, 0.80) & 0.0103 & 0.0091 & 0.010 & 80 \\
		\addlinespace
		& (1000, 0.40) & 0.0099 & 0.0088 & 0.010 & 81 \\
		& (1000, 0.50) & 0.0089 & 0.0084 & 0.011 & 78 \\
		& (1000, 0.60) & 0.0103 & 0.0069 & 0.008 & 87 \\
		& (1000, 0.70) & 0.0093 & 0.0069 & 0.009 & 80 \\
		& (1000, 0.80) & 0.0088 & 0.0063 & 0.010 & 75 \\
		\addlinespace
		\multicolumn{6}{@{}l}{\textit{Skewed-$t(3)$}} \\[2pt]
		& (250, 0.40) & 0.0201 & 0.0230 & 0.014 & 61 \\
		& (250, 0.50) & 0.0238 & 0.0244 & 0.012 & 71 \\
		& (250, 0.60) & 0.0276 & 0.0246 & 0.010 & 81 \\
		& (250, 0.70) & 0.0312 & 0.0244 & 0.007 & 71 \\
		& (250, 0.80) & 0.0353 & 0.0254 & 0.007 & 81 \\
		\addlinespace
		& (500, 0.40) & 0.0331 & 0.0289 & 0.007 & 84 \\
		& (500, 0.50) & 0.0178 & 0.0154 & 0.011 & 81 \\
		& (500, 0.60) & 0.0217 & 0.0161 & 0.008 & 81 \\
		& (500, 0.70) & 0.0250 & 0.0170 & 0.007 & 79 \\
		& (500, 0.80) & 0.0162 & 0.0110 & 0.010 & 80 \\
		\addlinespace
		& (1000, 0.40) & 0.0170 & 0.0150 & 0.010 & 79 \\
		& (1000, 0.50) & 0.0139 & 0.0097 & 0.010 & 82 \\
		& (1000, 0.60) & 0.0174 & 0.0117 & 0.008 & 90 \\
		& (1000, 0.70) & 0.0164 & 0.0095 & 0.009 & 80 \\
		& (1000, 0.80) & 0.0145 & 0.0074 & 0.010 & 77 \\
		\bottomrule
	\end{tabular}

	\begin{minipage}{\linewidth}\footnotesize
		\smallskip
		Calibration fraction controls a bias--variance trade-off in tail estimation: smaller $f_c$ reduces estimation variance but introduces bias due to insufficient tail coverage, while larger $f_c$ improves quantile estimation at the cost of reduced effective test sample size. For $T \leq 500$, no calibration fraction yields stable Basel Green performance under misspecification, reflecting the $O((\alpha T)^{-1/2})$ variance scaling implied by tail sparsity. For $T \geq 1{,}000$, performance stabilises for $f_c \in [0.60, 0.80]$, supporting the practical recommendation $f_c \geq 0.60$.
	\end{minipage}
\end{table}